\section{Physics Interpolation: Full-Device Velocity Fields}

\subsection{Architectural Role}

The interpolation module is a downstream consumer of full-device CFD velocity
fields. It does not generate geometry, regenerate meshes, solve Stokes systems,
or alter CFD outputs. The scientific workflow is:

\[
\mathrm{Hydraulics}\rightarrow \alpha(t),
\qquad
\mathrm{Full\text{-}device~CFD}\rightarrow \mathbf{u}(x,y;\alpha_i),
\qquad
\mathrm{Interpolation}\rightarrow \mathbf{u}(x,y;\alpha).
\]

Here \(\alpha\) is the left-branch flow fraction,

\[
  \alpha = \frac{Q_L}{Q_{\mathrm{in}}},
  \qquad
  Q_R = (1-\alpha)Q_{\mathrm{in}}.
\]

The present implementation interprets the two-dimensional CFD field as a
height-averaged background-flow approximation over the full loop domain. Future
ML modules can sample this continuous background field, but no transformer
integration is implemented in this module.

\subsection{Files}

\begin{itemize}
  \item \texttt{src/physics/interpolation/library.py}: discovers and validates
        the frozen CFD split library.
  \item \texttt{src/physics/interpolation/split\_interpolator.py}: chooses
        neighboring split cases and linearly interpolates P2 velocity
        coefficients.
  \item \texttt{src/physics/interpolation/field\_sampler.py}: evaluates an
        interpolated finite-element velocity field at arbitrary physical
        coordinates.
  \item \texttt{src/physics/interpolation/types.py}: typed data structures for
        loaded cases, interpolated fields, and sampled values.
  \item \texttt{src/physics/interpolation/validation.py}: withheld-case
        validation and scientific consistency checks.
  \item \texttt{scripts/physics/interpolation/validate\_velocity\_interpolation.py}:
        command-line validation entry point.
\end{itemize}

\subsection{Interpolation Model}

The frozen CFD library contains discrete solutions at

\[
  \alpha_i \in \{0.05,0.10,\ldots,0.95\}.
\]

For a requested \(\alpha\), the module finds neighboring stored fields
\(\alpha_{\mathrm{low}}\leq\alpha\leq\alpha_{\mathrm{high}}\) and computes

\[
  w =
  \frac{\alpha-\alpha_{\mathrm{low}}}
       {\alpha_{\mathrm{high}}-\alpha_{\mathrm{low}}}.
\]

Because all Version 1 CFD cases use the same production mesh, the same P2
velocity basis, and the same DOF ordering, interpolation is performed directly
on the finite-element velocity coefficients:

\[
  \mathbf{U}(\alpha)
  =
  (1-w)\mathbf{U}(\alpha_{\mathrm{low}})
  +
  w\mathbf{U}(\alpha_{\mathrm{high}}).
\]

This is coefficient interpolation, not interpolation of rendered figures,
streamlines, pressure plots, or regular-grid visualizations. If \(\alpha\)
matches a stored library fraction to the documented tolerance, the stored field
is returned exactly and no blend is performed.

\subsection{Sampling}

The public sampling operations name the coordinate frame explicitly:

\begin{verbatim}
library = VelocityFieldLibrary.from_directory(
    "outputs/physics/full_device_cfd/solutions"
)
field = library.interpolate(left_fraction=0.33)
samples_cfd = field.sample_cfd(points_cfd_um)
samples_device = field.sample_device(points_device_um, coordinate_convention)
\end{verbatim}

\texttt{sample\_cfd} is the low-level native operation. It expects points in
the frozen CFD frame \((x_{\mathrm{cfd}},y_{\mathrm{cfd}})\), where \(+y\)
points downward, and returns velocity components in the same frame.
\texttt{sample\_device} expects points in the device Cartesian frame, where
\(+y\) points upward, and returns device-frame velocity components. The older
\texttt{sample(points\_um)} method is retained only as a backward-compatible
alias for \texttt{sample\_cfd}.

The sampler reconstructs the scikit-fem P2 vector basis on the shared
production mesh and evaluates the FEM field directly:

\[
  \mathbf{u}(\mathbf{x};\alpha)
  =
  \sum_j U_j(\alpha)\boldsymbol{\phi}_j(\mathbf{x}).
\]

The output velocity units are meters per second in the requested output frame.
Points outside the CFD domain are flagged with \texttt{inside\_domain = false};
their velocity, speed, and direction values are returned as NaN. For
near-zero-speed points inside the domain, the velocity components are retained
while the unit direction is NaN.

\subsection{Validation}

The validation routine withholds stored CFD cases at
\(\alpha=0.20,0.40,0.60,0.80\). Each withheld field is reconstructed from its
two neighboring cases, then compared against the direct CFD result at the
shared velocity DOFs and at a common set of physical sample points inside the
junction domain.

\begin{center}
\begin{tabular}{rrrrrrr}
\toprule
\(\alpha\) & neighbors & \(w\) & DOF RMSE & sample speed RMSE & mean rel. error & mass residual \\
\midrule
0.20 & 0.15/0.25 & 0.500 & \(2.491{\times}10^{-15}\) & \(1.853{\times}10^{-15}\) & \(5.998{\times}10^{-14}\) & \(5.487{\times}10^{-17}\) \\
0.40 & 0.35/0.45 & 0.500 & \(1.951{\times}10^{-15}\) & \(1.955{\times}10^{-15}\) & \(4.282{\times}10^{-14}\) & \(6.754{\times}10^{-18}\) \\
0.60 & 0.55/0.65 & 0.500 & \(2.042{\times}10^{-15}\) & \(2.017{\times}10^{-15}\) & \(3.656{\times}10^{-14}\) & \(-4.136{\times}10^{-17}\) \\
0.80 & 0.75/0.85 & 0.500 & \(1.936{\times}10^{-15}\) & \(1.878{\times}10^{-15}\) & \(6.156{\times}10^{-14}\) & \(-8.948{\times}10^{-17}\) \\
\bottomrule
\end{tabular}
\end{center}

These near-machine-precision errors are expected because the current prescribed
split Stokes formulation is linear in the imposed outlet flux split on a fixed
mesh. The validation nevertheless remains important because it protects the
library contract: mesh identity, DOF ordering, field units, exact-match
behavior, and output immutability.

Scientific consistency checks confirmed:

\begin{itemize}
  \item no NaN values inside the valid sampled domain;
  \item maximum mass-balance residual
        \(1.135{\times}10^{-16}~\mathrm{m^2/s}\);
  \item inlet profile unchanged across tested \(\alpha\);
  \item maximum velocity on labeled no-slip wall facets equal to
        \(0.0~\mathrm{m/s}\);
  \item linear interpolation preserves the discrete incompressibility diagnostic
        to the same scale as the stored CFD solutions.
\end{itemize}

\begin{figure}[H]
  \centering
  \includegraphics[width=0.95\linewidth,height=0.75\textheight,keepaspectratio]{../outputs/physics/interpolation/validation/figures/withheld_error_summary.png}
  \caption{Withheld-case interpolation errors versus left-flow fraction. The errors are computed from FEM velocity coefficients and direct physical-domain samples, not from plotted images.}
  \label{fig:interpolation-withheld-summary}
\end{figure}

\begin{figure}[H]
  \centering
  \includegraphics[width=0.95\linewidth,height=0.75\textheight,keepaspectratio]{../outputs/physics/interpolation/validation/figures/withheld_0p40_velocity_errors.png}
  \caption{Representative withheld-case validation for \(\alpha=0.40\): interpolated speed, speed error, and angular error away from low-speed regions.}
  \label{fig:interpolation-withheld-040}
\end{figure}

\subsection{Execution}

\begin{verbatim}
.\.venv\Scripts\python.exe -m scripts.physics.interpolation.validate_velocity_interpolation --overwrite
\end{verbatim}

The validation outputs are written outside the frozen CFD library:

\begin{verbatim}
outputs/physics/interpolation/validation/
\end{verbatim}
